\documentclass{jss}

\author{
  Aydin Sareff-Hibbert\\The University of Sydney
  \And
  Gillian Z. Heller\\The University of Sydney
}

\title{\pkg{gamlss.longitudinal}: Longitudinal GAMLSS Models With Copula Dependence in \proglang{R}}

\Plainauthor{Aydin Sareff-Hibbert, Gillian Z. Heller}
\Plaintitle{gamlss.longitudinal: Longitudinal GAMLSS Models With Copula Dependence in R}
\Shorttitle{\pkg{gamlss.longitudinal}: Longitudinal GAMLSS Models}

\Abstract{
Longitudinal analyses often require flexible marginal distributions and
interpretable within-subject dependence, especially when scale, skewness, tails,
or association vary with time or covariates. The \proglang{R} package
\pkg{gamlss.longitudinal} implements first-order copula dependence for GAMLSS
margins in a formula-based workflow. Any supported marginal parameter and
copula parameter can depend on fixed or smooth covariates, allowing users to
model location, scale, shape, and adjacent-time dependence in one fitted object.
The package provides model fitting, distribution and copula screening,
inference, diagnostics, prediction, simulation, and reviewer-facing replication
workflows. A reproducible worked example and targeted validation studies
illustrate parameter recovery, diagnostic checks, optimizer tradeoffs, and the
relationship to existing \proglang{R} tools for GAMLSS, GEE/GLMM, bivariate
copula regression, and vine copulas.
}

\Keywords{GAMLSS, longitudinal data, copula, distributional regression, \proglang{R}}
\Plainkeywords{GAMLSS, longitudinal data, copula, distributional regression, R}

\Address{
  Aydin Sareff-Hibbert\\
  Faculty of Medicine and Health\\
  The University of Sydney\\
  NSW, Australia\\
  E-mail: \email{aydin.hibbert@sydney.edu.au}
}

\begin{document}

\section{Introduction}

% Replace this section using paper/jss-manuscript-blueprint.md.
% Center the software contribution and the gap relative to gamlss, gamlss2,
% geepack, lme4, GJRM, VineCopula, and rvinecopulib.

\section{Statistical model}

\subsection{Longitudinal GAMLSS margins}

% Define subject, time, response, marginal parameters, links, fixed terms, and
% smooth terms.

\subsection{First-order copula dependence}

% Define adjacent-pair copula dependence and state the exact first-order model
% assumptions. Be explicit about what is and is not implied for higher-order
% dependence.

\subsection{Discrete responses and incomplete panels}

% Explain rectangle probabilities and missing-panel contribution rules.

\subsection{Inference and uncertainty}

% Summarize covariance, robust/simulation guidance, and validation evidence.
% Move Hessian derivations to supplementary material.

\section{Software design and user interface}

\subsection{Package design}

% Explain gamlss_longitudinal(), family/copula metadata, fitted-object class,
% and stored model information.

\subsection{Minimal fit}

\begin{CodeChunk}
\begin{CodeInput}
R> fit <- gamlss_longitudinal(
+    mu.formula = response ~ time + group + s(baseline),
+    sigma.formula = ~ time + group,
+    theta.formula = ~ time + group,
+    dataset = dat,
+    time_var = "time",
+    subject_var = "id",
+    margin_dist = GA(),
+    copula_dist = "t"
+  )
R> summary(fit)
\end{CodeInput}
\end{CodeChunk}

\subsection{Selection helpers}

% Demonstrate select_margin(), select_copula(), and
% select_joint_distribution().

\subsection{Standard methods}

% Demonstrate summary(), coef(), vcov(), confint(), predict(), simulate(),
% plot(), logLik(), formula(), terms(), nobs(), fitted(), residuals(), and
% model.frame().

\subsection{Diagnostics}

% Show marginal diagnostics, copula diagnostics, Rosenblatt checks, tail checks,
% residual lag dependence, and missingness diagnostics.

\subsection{Optimizer controls}

% State when to use separate RS, joint RS, and CG, and how convergence metadata
% should be interpreted.

\section{Reproducible worked workflow}

% Use a public, package-shipped, or fully simulated primary example.
% Required steps: load/generate data, screen margins, screen copulas, fit,
% summarize, diagnose, predict, simulate, and interpret.

\section{Relationship to existing R software}

% Add the compact comparison table described in paper/jss-manuscript-blueprint.md.

\section{Validation and performance}

% Keep targeted validation only: continuous recovery, discrete recovery,
% optimizer guidance, first-order misspecification, and copula misspecification.
% Move large grids to supplementary material.

\section{Reproducibility}

% Document paper/replicate.R, smoke and expanded profiles, manifest, session
% info, output hashes, and private-data policy.

\section{Discussion}

% State intended use, limitations, and future work.

\bibliography{references}

\end{document}
